

% \rw{After the setup, Section~\ref{subsec:eval-base} evaluates all real service compositions, Sections~\ref{subsec:eval-q1} and~\ref{subsec:eval-q2} answer Questions~\#1 and~\#2, and Section~\ref{subsec:eval-calibration} relaxes the zero-shot setting.}

This section evaluates the performance of composite service models on unseen compositions and configurations, and thus, answers the two research questions. We complement this through experiments that show how to infer and optimize performance guarantees for different applications and SLOs.

\subsection{Experimental Setup}
%
To answer the two research questions (cf. Section~\ref{sec:motivation}), we require representative service compositions, where each service can be configured individually; a testbed for benchmarking individual services, as well as deploying the full service compositions; and lastly, baselines to evaluate the utility of our approach, according to key performance metrics. 

% \TODO{Must argue why these services are realistic and that they are built upon existing functionalities.}

\vspace{3pt}
\noindent
\textbf{Service Generation.} Customers can request applications from different domains (e.g., audio or lidar processing), so we evaluate \textit{Pegasus} across three of them. Every service builds on an abstract service interface~\cite{sedlak_multi-dimensional_2026} that provides queuing, monitoring, and multi-worker processing, and exposes the resources and software knobs. An LLM only inserts the target functionality into this wire-frame: given a one-sentence prompt, a reference implementation, and the desired knobs, it completes the functions that initialize a worker and process an item; Appendix~\ref{apx:service-generator} provides further details on the generator.

\vspace{3pt}
\noindent
\textbf{Applications.}
Each of the three \textit{application categories} (Audio, Lidar, Vision) has its own pool of 7--8 services, from which we build five compositions of 3--6 services (Table~\ref{tab:applications}). Every service has a counterpart in a published stream processing pipeline: video analytics~\cite{zhang_live_2017,jiang_chameleon_2018,romero_llama_2021}, point-cloud perception for autonomous driving~\cite{kato_autoware_2018,rusu_3d_2011}, or a speech       recognition front end~\cite{povey_kaldi_2011} (Table~\ref{tab:apx-services}). The expensive services run pretrained models (YOLOv10, YuNet, Vosk) on real recordings (KITTI, LibriSpeech). All services in a composition share \textit{one global data quality knob}; for example, every Vision composition uses the exact same frame height (up to 640 px).

\begin{table}[t]
    \centering
    \caption{Applications and their service compositions. Compositions share one data quality knob.
    % Vision also varies the model size (1--2) of the object and face detectors.
    \textbf{Bold}: the service with the largest share of the processing time. $K$: number of services; $P_{95}$: measured end-to-end $P_{95}$ latency of a run in ms, median over all runs (min--max). Appendix~\ref{apx:services} describes every service. Cov.: share of held-out runs below the pessimistic mean.}
    \label{tab:applications}
    \footnotesize
    \setlength{\tabcolsep}{3pt}
    \begin{tabular}{@{}l l c | l c@{}}
        \toprule
        {ID} & {Composition} & {$K$} & {$P_{95}$ [ms]} & {Cov.} \\
        \midrule
        \multicolumn{5}{@{}l}{\textit{Audio}; data quality: clip length, 2--12\,seconds} \\
        A1 & resample$\to$\textbf{denoise}$\to$MFCC & 3 & 70 (20--160) & .88 \\
        A2 & \textbf{denoise}$\to$voicedetect$\to$encode & 3 & 90 (30--180) & .95 \\
        A3 & resample$\to$voicedetect$\to$\textbf{pitch}$\to$encode & 4 & 90 (30--190) & \hl{.92} \\
        A4 & denoise$\to$voicedetect$\to$\textbf{transcribe} & 3 & 2110 (910--5480) & .97 \\
        A5 & resample$\to$pitch$\to$MFCC$\to$\textbf{transc.}$\to$encode & 5 & 2020 (750--4410) & .99 \\
        \midrule
        \multicolumn{5}{@{}l}{\textit{Lidar}; data quality: max.\ KITTI distance, 6--40\,meter} \\
        L1 & voxel$\to$\textbf{ground}$\to$cluster & 3 & 150 (30--200) & \hl{.94} \\
        L2 & \textbf{ground}$\to$normals$\to$encode & 3 & 150 (30--250) & .93 \\
        L3 & voxel$\to$\textbf{register}$\to$birdview & 3 & 100 (10--210) & \hl{.80} \\
        L4 & \textbf{ground}$\to$cluster$\to$birdview$\to$encode & 4 & 160 (40--230) & \hl{.84} \\
        L5 & voxel$\to$\textbf{normals}$\to$register$\to$cluster$\to$encode & 5 & 180 (20--340) & \hl{.88} \\
        \midrule
        \multicolumn{5}{@{}l}{\textit{Vision}; data quality: video frame height, 160--640\,pixel} \\
        V1 & \textbf{object}$\to$transform$\to$compress & 3 & 240 (30--1110) & \hl{.94} \\
        V2 & \textbf{face}$\to$transform$\to$edge & 3 & 120 (20--780) & .90 \\
        V3 & denoise$\to$\textbf{object}$\to$face$\to$compress & 4 & 260 (30--1310) & .93 \\
        V4 & \textbf{transform}$\to$denoise$\to$motion$\to$QR$\to$edge & 5 & 100 (30--230) & \hl{.92} \\
        V5 & motion$\to$QR$\to$\textbf{object}$\to$face$\to$edge$\to$compress & 6 & 280 (40--1550) & .93 \\
        \bottomrule
    \end{tabular}
\end{table}



\vspace{3pt}
\noindent
\textbf{Testbed.} Services are executed on one host, equipped with an AMD Ryzen 9 7950X and 128\,GB RAM. Each service worker (i.e., a Python process) consumes by default up to 1 CPU core. To fulfill Assumption~\ref{asm:load-independence}, we must avoid any CPU throttling; we achieve this by slightly over-provision the resources and allocating an additional half core ($0.5$) for every service.

\vspace{3pt}
\noindent
\textbf{Training Data (Isolated Benchmarks).} \textit{Pegasus} and the baseline methods use isolated service benchmarks for training per-service performance models.
%(cf. Sec.~\ref{subsec:learning}). 
We collect 150 training runs per service, where each run applies a randomly sampled service configuration from the configuration space in Table~\ref{tab:applications}; each runs uses a utilization levels $\rho=[0.1,1.0]$, and a number of workers $c=[1,4]$.
Every benchmarking run lasts 1 minute, excluding an additional 10s period for startup, and reports the mean execution time and variance under the configuration.

\vspace{3pt}
\noindent
\textbf{Testing Data (Service Compositions).} 
%
Measuring the true end-to-end latency for a service composition requires deploying and wiring all its services. For each service composition from Table~\ref{tab:applications}, we measure 150 runs of 1 minute duration each. As mentioned, we use one global \textit{data quality} level across a composition, which we again draw from a random uniform distribution. For every application, we run an 8\,s preliminary probe to set the load so that the slowest service (i.e., longest execution time) is 50\% utilized ($\rho=0.5$); Experiment~\ref{subsec:ex-varying-load} will evaluate higher application load. As for the training data, we configure each service to have workers $c=[1,4]$.


\vspace{3pt}
\noindent
\textbf{Baselines.}
%
We compare \textit{Pegasus} against five baseline methods. Like \textit{Pegasus}, each baseline is trained only on observations (or benchmarks) of individual services, never on a run of the evaluated composition $S$. The ground truth is the measured $(1-\epsilon)$ quantile of the end-to-end latency across $S$.
%
\begin{itemize}
    \item \textbf{No queueing (\textit{GP Sum}, \textit{DNN Sum}):} per service, a model predicts the
    execution time from the configuration; the end-to-end latency is the $(1-\epsilon)$ quantile of their sum.
    \textit{GP Sum} reuses the \textit{Pegasus}'s GPs, whereas \textit{DNN Sum} uses a feed-forward network with the standard deviation of its training residuals.
    GP and neural-network regressors are common performance models of configurable systems~\cite{alipourfard_cherrypick_2017, van_aken_automatic_2017, li_metis_2018, fu_use_2021,liang_modular_2023}; summing their per-service predictions is the most direct way to compose their performance without data of the composition.

    \item \textbf{Queueing formula (\textit{GP Queue}):} the same GP predictions enter the Allen--Cunneen approximation of the mean waiting time in a G/G/$c$ queue~\cite{allen_probability_1990}, which reduces to Kingman's formula~\cite{kingman_single_1961} for one worker, with an exponential waiting-time tail~\cite{harchol2013performance}; the per-service $(1-\epsilon)$ quantiles are added. \textit{GP Queue} accounts for queueing, but it is an approximation, not a bound, and it uses GP predictions without a margin for epistemic uncertainty.

    \item \textbf{Learned quantiles (\textit{Quantile Union}, \textit{Conformal}):}
    the same GPs are used to predict every service's $(1-\epsilon/K)$ latency quantile;
    in \textit{Quantile Union}, the $K$ predictions are added as a union bound. \textit{Conformal} additionally multiplies each prediction by a split-conformal correction factor calibrated on held-out training configurations~\cite{vovk_conformal_2020}. In contrast to \textit{Pegasus}, both baselines learn the waiting time from the benchmarks instead of deriving it from the arrival process; their predictions therefore hold only for arrivals like those in the benchmarks (Poisson), and a different arrival pattern requires new benchmarks.
\end{itemize}

As a reference, an \textit{Oracle} adds the \textit{measured} $(1-\epsilon/K)$ latency quantiles of the services inside the composition. The \textit{Oracle} uses the ground truth and is not a competitor; it shows the tightest result any union bound over services can reach.

\vspace{5pt}
\noindent
\textbf{Performance Metrics.}
%
We evaluate the performance of \textit{Pegasus} and the baseline methods according to three key metrics: 
%
\begin{itemize}
    \item \textit{\textbf{tightness}} of the latency bound, which we compute as the predicted $P_{95}$ latency divided by the measured $P_{95}$. The \textit{tightness} should have a slight slack; below $1\times$, the SLO is violated, while above, resources are wasted. 
    %
    \item \textit{\textbf{coverage}} of the latency bound, which we calculate as the number of runs across all seeds that are not below the 95\% confidence interval of its measured $P_{95}$. A method holds if, under every seed, it covers at least $\gamma = 95\%$ of the runs; below, we say the bound does not hold.
    %
    \item \textit{\textbf{complexity}} of providing the latency bound, which we measure through two splits: the time required for training the per-service performance models, and the time for inferring a latency bound for a target composition.
    % to evaluate all held-out runs of one seed. Pegasus, \textit{GP Sum}, and \textit{GP Queue} share the same GPs and thus the training time.
\end{itemize}


\subsection{Base Case}

% \TODO{Decide on experiment structure}

% \subsubsection{Base Case}
\label{subsec:eval-base}

We analyze the tightness and coverage of end-to-end latency bounds across different methods, for which we use the 15 application instances presented. We randomly draw 50 benchmarking runs per service (with \hl{ten} seeds), which presents the available training data; \textit{Pegasus}, among other methods, uses this data for fitting the performance predictor (the GP). For every service in a held-out run, we also record whether its measured mean execution time lies below the pessimistic mean $m_i^{+}$ at $\gamma = 0.95$, which shows whether $\gamma$ is calibrated.

\paragraph{Result}
%
\textbf{\textit{Pegasus} guarantees hold on all held-out compositions and configurations.}
%
Fig.~\ref{fig:results-b1} shows that \textit{Pegasus} holds on all 15 applications at a median tightness of \hl{$1.97\times$}; \textit{Quantile Union} holds on \hl{5} (\hl{$1.46\times$}). \textit{GP Sum} and \textit{DNN Sum} fail to provide a guarantee for any chains, since they ignore queueing; 
\textit{GP Queue} ignores the epistemic uncertainty and holds on only \hl{3} of 15 applications (\hl{1.17}$\times$). \textit{Conformal} holds on all at \hl{4.94}$\times$.
%
\textbf{$\mathbf{\gamma}$ is slightly overstated per service.} On held-out benchmark runs, the pessimistic mean covers 0.92 of the runs at $\gamma = 0.95$, and inside the compositions \hl{0.92} of the services (Cov. in Table~\ref{tab:applications}; lowest for L3 and L4). 
% All services of a run are covered in only 72\% of the runs, below the 82\% that independent errors would give, so the errors of the services in a run are positively correlated.
%
% To decrease the clutter from reporting non-competitive methods, we do not list \textit{DNN Sum} in the remaining experiments.

\begin{figure}[t]
    \centering
    \includegraphics[width=1.0\linewidth]{figures/experiments/b1_all_chains_benchmark.pdf}
    \caption{\textit{Pegasus} gives performance guarantees that hold for all evaluated service compositions; compare to \textit{Conformal}, the only other guarantee among baselines, \textit{Pegasus} holds on average $7\times$ more accurate. Evaluated applications span three categories (audio, lidar, vision) and compose 22 different services into chains of 3--6 sequential services. Appendix~\ref{apx:scenarios} holds the service details.}
    \label{fig:results-b1}
\end{figure}

% \PROBLEM{The slack of $1.72\times$ is unexplained, and it is unclear    
%   which variant                                                           
%   ``Pegasus'' is.                                                         
%   (1) Slack decomposition: the Oracle is close to $1\times$, so the union 
%   bound itself is                                                         
%   cheap; most slack must come from the epistemic shift, the               
%   Chernoff/Kingman tail, or the                                           
%   cyclic-assignment pessimism. An ablation that removes these one at a    
%   time would explain the                                                  
%   gap and answer the obvious question: why not Quantile Union             
%   ($1.6\times$) plus a small margin?                                      
%   (2) Node vs.\ SNC: Sec.~IV defines two variants, but Figs.~3--7 show a  
%   single ``Pegasus''                                                      
%   line, and only Sec.~V-E2 reports both. State for every figure which     
%   variant it shows, and back                                              
%   the ``when to choose which'' criteria with data.}

\subsection{Model Training and Bound Calibration}

\subsubsection{Varying Evidence per Service}

Benchmarking is costly, therefore we analyze the accuracy gain of additional training runs on the service performance models, which is expected to improve the tightness and coverage of \textit{Pegasus} guarantees. We use the \textit{V3} application from Table~\ref{tab:applications} and train each service's GP with 3 to 200 runs; these models are then used to bound the performance of held-out runs. By default, these training and testing splits contains benchmarks with various worker and load levels; samples are drawn according to \hl{10} seeds.

\paragraph{Result}
\textbf{Twenty benchmarks per service suffice to obtain a reliable guarantee. }
%
Fig.~\ref{fig:results-b2a} shows the prediction errors on service execution time for each service in \textit{V3}, whereas Fig.~\ref{fig:results-b2b} shows the respective \textit{Pegasus} guarantees when bounding the performance according to these models; additionally, Fig.~\ref{fig:results-gp-denoiser} visualizes the GP for one trained service (\textit{denoiser}). With increasing benchmark runs, two factors set in: (1) the prediction models become increasingly accurate; all service errors converge to $3\%$ to $15\%$, and the \textit{predicted means} gives the absolute difference between the GP's prediction (without uncertainty) and the measured latencies. 
% \tmi{The impact of the services' prediction errors here is highly disproportional: the notable gain on coverage is caused by the slowest service (object detector) is the main reason why Pegasus converges into a stable guarantee.}
Also, (2) the GP's uncertainty decreases as the configuration space gets increasingly covered.
%
At 3 and 5 runs, the \textit{Pegasus} bounds are infinite due to this epistemic uncertainty; after 10 runs \textit{Pegasus} produced finite guarantees, though with high uncertainty in some seeds, and at 50 runs we get stabler results (no seed over $3\times$).


\begin{figure}[t]
    \centering
    \subfloat[Execution time and uncertainty]
    {\includegraphics[height=0.531\linewidth]{figures/experiments/b2_evidence_real_uncertainty_execution.pdf}
    \label{fig:results-b2a}}
    \hfill
    \subfloat[Impact on bound per service]
    {\includegraphics[height=0.531\linewidth]{figures/experiments/b2_evidence_real_uncertainty_cost.pdf}
    \label{fig:results-b2b}}
    \caption{\textit{Pegasus} requires roughly 20 individual benchmarks per service to provide an accurate guarantee for the \textit{V3} composition. Left: the prediction error gradually declines and converges with more training runs per service. The \textit{detector} service contributes the most ($\sim 76\;$ms) to the latency. Right: initially, the \textit{Pegasus} bound is loose ($2.5\times$ at 10 runs), owed to the high prediction error in the other services, especially the \textit{detector}. The \textit{precited means} are extracted from the same GP, but without uncertainty.}
    \label{fig:results-b2}
\end{figure}

\vspace{5pt}
\subsubsection{Calibrating Bounds with End-to-End Samples}
\label{subsubsec:fine-tuning}

% \TODO{Build C1 upon the end-to-end benchmarks from C2}

Service providers need to give guarantees for compositions prior to deploying them, which is what \textit{Pegasus} addressed; however, as soon as end-to-end evidence becomes available, it should be possible to integrate it into the guarantee, ideally tightening the bound. To that extent, we split the 150 runs of the \textit{V1} composition into $n = 100$ end-to-end samples and 50 test runs. \textit{Pegasus} multiplies its bound by the $\lceil (n+1)\gamma \rceil$-th smallest ratio of measured to predicted $P_{95}$ among the samples (split-conformal), which needs $n \geq 20$. We construct two \textit{Pegasus} guarantees: \textit{pretrained} from the existing GPs, or \textit{blank}, without any knowledge (i.e., uses 50 end-to-end samples to train the GP, and the rest for calibration). The blank version thus has the same data as the baseline: 
% \textit{Quantile Sum} adds the per-service $P_{95}$ measured inside the chain, and 
in \textit{Fluxion}~\cite{liang_modular_2023}, the authors chained one GP per service that takes the service configuration and the predicted quantiles of the preceding service, all done on end-to-end samples.

\paragraph{Result}
\textbf{End-to-end samples tighten Pegasus without breaking it.} Fig.~\ref{fig:results-c1} shows that \textit{Pegasus (pretrained)} holds for every $n$.
% From $n = 20$, the calibration takes effect: at $n = 50$, Pegasus tightens from $1.58\times$ to $1.35\times$ and Pegasus SNC from $2.59\times$ to $1.38\times$. \textit{Quantile Sum} never holds, the \textit{Fluxion}-style model only at $n = 50$.
From $n = 20$, the calibration takes effect: at $n = 100$, it tightens from \hl{$1.83\times$} to \hl{$1.37\times$}. \textit{Pegasus (blank)}, trained on the samples alone, holds at all $n$ but \hl{two ($n = 20$ and $n = 80$)} and reaches \hl{$1.29\times$} at $n = 100$. Neither \textit{Quantile Sum} nor the \textit{Fluxion}-style model holds at any $n$.


% \begin{figure}[t]
%     \centering
%     \subfloat[abcd]{
%     \includegraphics[width=0.485\linewidth]{figures/experiments/c1_finetuning_real.pdf}}
%     \hfill
%     \subfloat[GP for \textit{transformer} service]{
%     \includegraphics[width=0.485\linewidth]{figures/transformer_gp_quality_load.pdf}}
%     \caption{Pegasus can improve its tightness on the \textit{V1} composition by correcting the slack according to a split conformal factor; this uses up to 100 end-to-end samples. \textit{Fluxion} requires end-to-end samples, which it composes into a performance predictor. (\textbf{Right figure must be one of the services in it; alternatively, split into two figures})}
%     \label{fig:results-c1}
% \end{figure}

\begin{figure}[t]
    \centering
    \begin{minipage}[t]{0.485\linewidth}
        \centering
        \includegraphics[width=\linewidth]{figures/denoiser_gp_quality_load.pdf}
        \caption{GP for \textit{denoiser} service. \hl{how many samples used here?}}
        \label{fig:results-gp-denoiser}
    \end{minipage}
    \hfill
    \begin{minipage}[t]{0.485\linewidth}
        \centering
        \includegraphics[width=\linewidth]{figures/experiments/c1_finetuning_real.pdf}
        \caption{\textit{Pegasus} improves its tightness on the by correcting the slack according to a split conformal factor, using up to 100 end-to-end samples of \textit{V1}.}
        \label{fig:results-c1}
    \end{minipage}
\end{figure}

\subsection{Extrapolating Configurations (Busier, Larger, and Better)}

\subsubsection{Varying Load and Arrival Pattern}
\label{subsec:ex-varying-load}

When and how a queue is formed is mainly determined by the arrival process---more specifically by the number of arrivals and their spacing. Performance guarantees must hold across these arrival processes; we analyze this in detail for the $V1$ application and perform two separate experiments where we (1) vary the arrival rate $\rho$ between $[0.1,1.0]$ (i.e., from low to complete saturation) for a Poisson arrival process, and (2) at a stable $\rho = 0.5$, vary the variation of gamma-distributed gaps ($C_a^2$) between $[0,4]$.  \rw{The execution-time GPs stay unchanged; Pegasus receives the arrival process and \textit{GP Queue} its rate and $C_a^2$. \textit{Quantile Union} and \textit{Conformal} are retrained for every setting.} 

\paragraph{Result}
%
\textbf{Pegasus holds even close to service saturation. }
%
Fig.~\ref{fig:results-b4a} shows that Pegasus holds at every load, \hl{1.62--1.91}$\times$ up to $u=0.7$, \hl{3.89}$\times$ at $u=0.9$ (\hl{38\%} infinite bounds).

\textbf{Pegasus guarantees hold across various arrivals. } 
%
Fig.~\ref{fig:results-b4b} shows that Pegasus holds for every $C_a^2$, from \hl{1.13}$\times$ ($C_a^2=0$) to \hl{2.18}$\times$ ($C_a^2=4$). Quantile Union holds up to \hl{$C_a^2=0.3$}, Conformal up to $C_a^2=2$; GP Queue never holds.

\begin{figure}[t]
    \centering
    \subfloat[Low to high service utilization]
    {\includegraphics[height=0.531\linewidth]{figures/experiments/b4a_load_real.pdf}
    \label{fig:results-b4a}}
    \hfill
    \subfloat[Periodic to cluttered arrivals]
    {\includegraphics[height=0.531\linewidth]{figures/experiments/b4b_arrivals_real.pdf}
    \label{fig:results-b4b}}
    \caption{\textit{Pegasus} holds its performance guarantees across different arrival processes; here evaluated for the \textit{V1} composition. Left: the utilization of \textit{V1} is gradually increased, until its slowest service (\textit{detector}) becomes unstable and its queue grows infinite. Right: the variability in the arrival process is gradually increased; $0 \leq C_a^2 < 1$ define periodic arrivals with increasing jitter, $C_a^2=1$ is Poisson arrival, and beyond $C_a^2 > 1$, arrivals are increasingly clustered. }
    \label{fig:results-b4}
\end{figure}


\subsubsection{Varying Application Depth}
\label{subsec:ex-chain-depth}

While the 15 applications in the base case already spanned chains of 3--6 services, this experiment analyzes the scalability on larger, non-linear applications.
%
To that extent, we create increasingly larger map-reduce applications from the Vision services in Table~\ref{tab:applications}: 
% We construct the application as a hierarchical, balanced tree consisting of
\textit{object} and \textit{face detector} as leaf nodes, and transformers as parent nodes, performing the fan-in operations on processed video frames. The application size grows exponentially; at height $h=6$, it contains a total number of $(2^6 -1) = 63$ services.
% , with $16$ of them being detectors at the leave nodes.
We set utilization to $\rho=0.5$ of the slowest service (object detector), use $c=1$ worker per service, and capture 100 runs across randomly drawn configurations. 
Each method bounds all leaf-to-root paths (i.e., up to 32 paths for the leave nodes at $h=6$) and reports the maximum expected latency.

\paragraph{Result}
%
\textbf{Pegasus holds on larger applications, but looses tightness. }
%
slightly looses tightness with $1.3\times$ at every additional level.
%
Most importantly, the host resources (16 cores) did not suffice to execute this application size (up to 63 services). At medium load ($\rho=0.5$), all services exhaust half of their allocated resources at most; still, after $h=3$ the host's resources are exhausted and the services start throttling. This slows down service execution and increases the error of prediction models; Fig.~\ref{fig:results-b3b} shows how the services in the map-reduce application become increasingly slower as the overall number of actively executed services on the host grows; \textit{face detector} operates up to 55\% slower than expected at height $h=6$. Consequentially, all methods loose tightness, with \hl{\textit{Pegasus} being the only method whose slack could compensate the prediction errors at $h=6$ (\textit{Conformal} no longer holds there)}.

\textbf{Pegasus gives bounds within seconds. }
%
To analyze \textit{Pegasus} tightness and complexity despite this hardware contention, we simulate the remaining trees' performance using the unconstrained performance records. As shown in Fig.~\ref{fig:results-b3-complexity-a}, training the 61 GPs at $h=6$ takes less than $10\;\text{s}$. Afterwards, we use these models to infer the performance guarantees, which takes less than $1\;\text{s}$ at $h=6$. Combined, this shows that Pegasus can train and infer a performance bound for 63 services in roughly 10 seconds. Note that the map-reduce application repeats three services, so it would be logical to cache and reuse GPs; we did not apply this here to analyze worst-case complexity.

\begin{figure}[t]
    \centering
    \subfloat[abcd]{\includegraphics[width=0.64\linewidth]{figures/experiments/b3_trees_real_large_tightness.pdf}
    \label{fig:results-b3a}}
    \hfill
    \subfloat[abcd]{\includegraphics[width=0.33\linewidth]{figures/experiments/b3_slowdown_concurrency.pdf}
    \label{fig:results-b3b}}
    \caption{All methods, including \textit{Pegasus}, slightly loose tightness with increasing services in a hierarchical map-reduce application. Left: From $h=4$ (15 services) on, the services start to throttle because the resources are saturated. Consequentially, the methods become inaccurate because the predicted execution time deviates increasingly from the actual execution time; in comparison to the other methods, \textit{Pegasus}' slack partially absorbs this error. Right: Services increasingly slow down as the hardware is saturated. Consequentially, the Pegasus guarantees become inaccurate.}
    \label{fig:results-b3}
\end{figure}

\begin{figure}[t]
    \centering
    \subfloat[Training and calibration time.]
    {\includegraphics[height=0.531\linewidth]{figures/experiments/b3_trees_real_large_training.pdf}
    \label{fig:results-b3-complexity-a}}
    \hfill
    \subfloat[Performance inference time.]
    {\includegraphics[height=0.531\linewidth]{figures/experiments/b3_trees_real_large_evaluation.pdf}
    \label{fig:results-b3-complexity-b}}
    \caption{Complexity for training and inference across methods in the hierarchical map-reduce application. Left: Pegasus shares its training time with other methods that use a GP as prediction model; the training time plateaus after $h=2$ because new nodes are always repetitions of existing services (\textit{object}, \textit{face}, \textit{transformer}). Right: the inference time; \textit{Pegasus} infers performance guarantees for a tree with $h=6$ (63 services) in roughly \hl{$570\;$ms}.}
    \label{fig:results-b3-complexity}
\end{figure}

\subsubsection{Custom Service Composition and Configuration}

\begin{table}[t]
    \centering
    \caption{Three selected evaluation scenarios, containing functional requirement, hard SLO thresholds, and optimization targets. Below: the selected composition and configuration.}
    \label{tab:scenarios}
    \footnotesize
    \setlength{\tabcolsep}{3pt}
    \begin{tabular}{llll}
        \toprule
        & \#1 (Audio) & \#10 (Lidar) & \#13 (Vision) \\
        \midrule
        Client & shared gateway & robotaxi & forensic review \\
        Rate [per s] & 6 & 20 & 12 \\
        Requires & \textit{denoise} & \textit{encode} & \textit{object} detector \\
        Hard SLOs & $c \leq 4$ workers & $P_{95} \leq 0.2$s, $c \leq 10$ & $q \geq 512$px, model 2 \\
        Optimized & qu. (1), lat. (0.5) & quality (1) & costs (1), lat. (0.1) \\
        \midrule
        Plan: chain & A1 & L4 & V1 \\
        quality & 12\,s & 40\,m & 512\,px \\
        workers & 1, 2, 1 ($c=4$) & 4, 2, 1, 2 ($c=9$) & 3, 1, 1 ($c=5$) \\
        bound [ms] & 126 & 198 & 1265 \\
        \bottomrule
    \end{tabular}
\end{table}

Our processing services can be composed and configured for optimizing customer experience. We already evaluated \textit{Pegasus} performance on held-out runs (compositions and configurations), so the remaining piece is to infer deployments from client requirements and verify whether they hold in practice.
%
We evaluate 15 scenarios that provide different SLOs (i.e., thresholds for \textit{latency}, \textit{data quality}, and \textit{costs}) and functional requirements (e.g., must include \textit{detector} service); Table~\ref{tab:scenarios} presents the details of three selected scenarios (client requirements and plan that \textit{Pegasus} infers), Appendix~\ref{apx:scenarios} holds the details and results for all of them.


For each scenario, we (1) filter the available service chains according to their provided functionality, and (2) filter their configurations according to predicted SLO fulfillment; lastly, from the remaining service configurations, we (3) select the configuration that optimizes the remaining, unconstrained SLO dimensions.
%
For example, in Scenario $\# 10$ the client requested an application including an \textit{encode} service, latency $P_{95} \leq 200\; \text{ms}$, and resource costs $c \leq 10\;\text{workers}$. The unconstrained dimension is \textit{quality}, so we select the configuration with the highest \textit{quality} that fulfills all hard requirements.
%
% \tmi{If multiple dimensions are unconstrained, the targets are optimized according to a weight, e.g., in $\# 3$ the client specified only a quality $q \geq 10\; \text{sec}$, leaving the other two dimensions for optimization with a clear focus on low cost ($w=1$) and less important latency (weight $w=0.1$).}
%
The resulting applications are deployed with the specified service and resource configurations and we capture 50 runs per scenario.


\paragraph{Result}

\textbf{Pegasus guarantees hold on the testbed.} Fig.~\ref{fig:results-c2} shows the predicted $P_{95}$ bound of the candidate compositions over the data quality $q$ (star: selected plan) and the measured $P_{95}$ of the 50 deployed runs (magenta bar). In all three scenarios, the bound lies above the measured $P_{95}$, at a tightness of roughly $1.4$--$1.8\times$, and the plan sits where an SLO binds: the worker limit (\#1), the latency limit (\#10), and the quality limit (\#13).




\begin{figure}[t]
    \centering
    \subfloat[Scenario \#1 (Audio)]{\includegraphics[width=0.318\linewidth]{figures/experiments/c2_tradeoff_s_a_1.pdf}
    \label{fig:results-c2a}}
    \hfill
    \subfloat[Scenario \#10 (Lidar)]{\includegraphics[width=0.318\linewidth]{figures/experiments/c2_tradeoff_s_l_5.pdf}
    \label{fig:results-c2l}}
    \hfill
    \subfloat[Scenario \#13 (Vision)]{\includegraphics[width=0.318\linewidth]{figures/experiments/c2_tradeoff_s_v_3.pdf}
    \label{fig:results-c2v}}
    \caption{Predicted $P_{95}$ bounds for three selected scenarios. Star: the selected plan (composition and configuration) with $c$ workers; box plot and magenta bar: measured latencies and their $P_{95}$ over 50 runs; red zone: SLO violated. Each star lies above its magenta bar, so the guarantee holds.}
    \label{fig:results-c2}
\end{figure}